\documentclass{article}

\usepackage[T1]{fontenc}

\usepackage{booktabs,graphicx,siunitx}

\begin{document}

\section{Chi-square goodness-of-fit}

\begin{tabular}{llll}
\toprule
Category & Observed & Expected & Std. residual \\
\midrule
discussion & 16 & 20.00 & $-$1.26 \\
lecture & 13 & 12.00 & 0.35 \\
seminar & 11 & 8.00 & 1.19 \\
\bottomrule
\end{tabular}

\begin{tabular}{llll}
\toprule
$\chi$\textsuperscript{2} & df (k$-$1) & p & Cohen's w [95\% CI] \\
\midrule
2.01 & 2 & .366 & 0.22 [0.00, 2.00] \\
\bottomrule
\end{tabular}

\includegraphics{figures/01_chi-square-goodness-of-fit/figure_bar.png}

Tests whether one categorical variable's counts depart from an expected distribution (equal by default). A p below alpha means the observed split differs from expected; a standardized residual beyond about $\pm$1.96 flags a category that significantly drives the result. Valid only when expected counts are mostly $\ge$ 5 (some texts allow $\ge$ 1 if no more than 20\% are below 5); for sparse categories use an exact / simulated test instead. Your significance threshold ($\alpha$) is 0.05.

A goodness-of-fit test, $\chi$\textsuperscript{2}(2, N=40)=2.01, p = .366, w=.22 [.00, 2.00].

\end{document}